\par
\SourceLeaf{368}%
Geben drei auf einander folgende Scalen
\[
M,0,0,0
\]
und sind die betreffenden Drehungspunkte $P,P',P''$,\\
\hspace*{5em}die Winkel $\lambda,\lambda',\lambda''$,\\
so sind $\tfrac12\lambda,\tfrac12\lambda',\tfrac12\lambda''$ die Winkel des sphärischen Dreiecks zwischen $P,P',P''$.

\Needspace{14\baselineskip}
\begin{center}[II.]\end{center}
\begin{center}\bfseries MUTATIONEN DES RAUMES VON DREI DIMENSIONEN.\end{center}
\begin{center}1.\end{center}

Seit langer Zeit habe ich das Verfahren gebraucht, jede Mutation eines Raumes (worunter ich verstehe seine Verrückung in einem andern Raume verbunden mit einer allgemeinen Vergrösserung oder Verkleinerung des erstern, unbeschadet der Conservation der Ähnlichkeit) durch vier Grössen
\[
a,b,c,d
\]
auszudrücken, deren Complex ich die Mutationsscale nenne. Es wird zugleich angenommen, dass Ein Punkt des beweglichen Raumes in dem andern als fest oder absolut betrachteten immer unbewegt bleibe, welcher Punkt mit Nullpunkt bezeichnet werden mag.

\Needspace{9\baselineskip}
\begin{center}2.\end{center}

Die Bedeutung dieser vier Grössen ist hierbei folgende. In dem festen Raume werden drei von $0$ ausgehende, rechte Winkel mit einander machende gerade Linien angenommen $01,02,03$. Man setze
\[
\BareRoot{(bb+cc+dd)}=\rho,
\]
und lasse
\[
\frac b\rho,\ \frac c\rho,\ \frac d\rho
\]
